\scanpage{16}
130\hfill J.-P. \textsc{Serre}

ou encore
\[
v(x)\leq\psi(v(f(x))).
\]

De plus, pour $v(x)\inserted{<}c_2=c_1/(p^h-1)$, la valuation du terme correspondant à $n=p^h$ est strictement inférieure à celle des autres termes, et l'on a donc
\[
v(f(x))=p^hv(x).
\]

Soit $d_n$ la borne supérieure des $v(x)$, pour $x\in T'_n$. Les $d_n$ sont $>0$. Si $x\in T'_n$, on a $f(x)\in T'_{n-1}$; d'où:
\[
d_n\leq\psi(d_{n-1})\quad\text{et}\quad d_n\leq\psi^{(n-1)}(d_1).
\]

Mais, pour tout $\alpha>0$, les itérés $\psi^{(n)}(\alpha)$ tendent vers 0. Il existe donc un entier $n_0$ tel que, si $n\geq n_0$, on ait $d_n<c_2$, d'où $d_{n+1}=d_n/p^h$. On obtient ainsi l'existence d'une constante $c_3>0$ telle que:
\[
v(x)\leq c_3/p^{nh}\quad\text{si}\quad x\in T'_n.
\]

Soit $c_4$ la valuation d'une uniformisante de $K$, et soit $e$ l'indice de ramification de $K(x)/K$. On a évidemment:
\[
v(x)\geq c_4/e.
\]

En comparant, on trouve $e\geq c\cdot p^{nh}$, avec $c=c_4/c_3$, ce qui démontre le lemme.

On voit ainsi que la condition (a) du théorème 4 est automatiquement vérifiée en dimension 1. La condition (c) l'est aussi pourvu que $h\geq2$, puisque $n_1=1$ et $n_0=h-1$. Comme le cas $h=1$ est trivial, on en déduit finalement:

\textbf{Théorème 5.} -- \textit{Soit $F$ un groupe formel défini sur $A$, de dimension 1 et de hauteur finie $h$. Supposons que $\operatorname{End}(F)=\mathbf Z_p$. L'image $G$ de $\operatorname{Gal}_K$ dans $\operatorname{Aut}(V)$ est ouverte.}

(L'hypothèse $\operatorname{End}(F)=\mathbf Z_p$ signifie que $F$ n'a pas de «multiplication complexe formelle», au sens de \textsc{Lubin} [2].)

\textbf{Remarques.} -- 1) Lorsqu'on ne fait plus l'hypothèse que $\operatorname{End}(F)=\mathbf Z_p$, on peut quand même déterminer l'algèbre de Lie $\mathfrak g$ de $G$. On trouve que c'est le \textit{commutant} de $\operatorname{End}(F)$ dans $V$. La démonstration est analogue à celle du cas particulier traité ici, mais sensiblement plus compliquée.
\margin{voir Sen (\unsure{Annals})\quad(}

2) Les résultats ci-dessus s'appliquent au \textit{groupe d'inertie} $I$ de $G$; cela se voit en passant au corps $\widehat K_{nr}$, et en remarquant que $\operatorname{End}(F)$ ne change pas par cette opération, en vertu de \textsc{Lubin} [2], \S\ 2 (c'est là une propriété
\margin{15}
